%! AP Statistics — Unit 9, Lesson 9.2 — Warm-Up (Answer Key)
\documentclass[10pt]{article}
\usepackage{apstats-article}
\usepackage{apstats-key}

\begin{document}

\pageheader{Unit 9, Lesson 9.2}{Warm-Up --- Answer Key}
\namedateperiod

\vspace{0.15in}

A study of 15 randomly selected college students recorded the number of hours spent
studying per week ($x$) and the score on a standardized exam ($y$).
The LSRL is $\hat{y} = 42.3 + 3.8x$.

\begin{enumerate}

  \item \textbf{Interpret the slope} of the LSRL in context.

        \ans{For each additional hour studied per week, the predicted standardized exam
        score increases by approximately 3.8 points.}

  \item One student studied 10 hours per week and earned a score of 81.
        \textbf{Calculate the residual} for this student.

        \vspace{0.05in}
        $\hat{y} = 42.3 + 3.8($\ans{10}$) = $ \ans{80.3}

        \vspace{0.08in}
        Residual $= y - \hat{y} = $ \ans{81} $-$ \ans{80.3} $= $ \ans{0.7}

        \vspace{0.1in}

  \item In Lesson 9.1, you learned about the sampling distribution of the slope $b$.
        Explain what it would mean if the \textbf{true slope $\beta = 0$} for this
        population.

        \ans{A true slope of $\beta = 0$ would mean there is no linear relationship
        between hours studied and exam score in the population --- knowing how many hours
        a student studies provides no useful information for predicting their score.}

\end{enumerate}

\end{document}
